\section*{Capítulo \ref{ch:b2:speciation} --- Especiação e macroevolução}

\begin{solution}{exo:b2:speciation:1}
Uma espécie é um grupo de populações naturais que se cruzam entre si,
reprodutivamente isolado de outros grupos desse tipo. O conceito falha
para as linhagens assexuadas (bactérias, rotíferos bdeloides,
dentes-de-leão: conceito morfológico ou filogenético, agrupamentos de
formas semelhantes ou linhagens diagnosticáveis); para os fósseis
(conceito morfológico, pois nenhum cruzamento pode ser feito); e para as
populações alopátricas que nunca se encontram (conceito ecológico ou
filogenético, pois o intercruzamento não pode ser testado) --- e ele é
embaralhado pelos carvalhos que hibridizam e pelas espécies em anel.
\end{solution}

\begin{solution}{exo:b2:speciation:2}
Rãs em meses diferentes: pré-zigótico (temporal). Mula: pós-zigótico
(esterilidade híbrida). Espermatozoide que não consegue se ligar:
pré-zigótico (gamético). Híbridos vigorosos de \omterm{def:b2:angiosperm-reproduction:seed}{sementes} estéreis:
pós-zigótico (esterilidade híbrida, ou colapso híbrido se a falha
aparece na geração seguinte). Grilos de cantos diferentes: pré-zigótico
(comportamental).
\end{solution}

\begin{solution}{exo:b2:speciation:3}
Alopátrica: uma barreira divide uma área de distribuição (os
camarões-pistola dos dois lados do istmo do Panamá). Peripátrica: alguns
fundadores além da área de distribuição (os tentilhões de Darwin a
partir de um ancestral continental). Simpátrica: sem barreira (a
mosca-da-maçã sobre o espinheiro e a macieira; um trigo \omterm{def:b2:genome-diversification:polyploidy}{alopoliploide}).
Parapátrica: ao longo de um gradiente com uma \omterm{prop:b2:speciation:reinforcement}{zona híbrida} no meio
(gramíneas dentro e fora de rejeitos de mineração, que florescem em
épocas diferentes).
\end{solution}

\begin{solution}{exo:b2:speciation:4}
A espécie 1 sempre vence; a espécie 2 sempre vence; coexistência
estável; qualquer uma vence, conforme os efetivos iniciais. A
coexistência é estável quando cada espécie pode invadir a outra em sua
capacidade de suporte, isto é, quando cada espécie limita a si mesma
mais do que limita a outra ($K_1 > \alpha K_2$ e $K_2 > \beta K_1$).
\end{solution}

\begin{solution}{exo:b2:speciation:5}
$\alpha K_2 = 400 < K_1 = 1000$: a espécie 1 invade; $\beta K_1 = 600
< K_2 = 800$: a espécie 2 invade. $N_1^{*} = (1000 - 400)/(1 - 0.3) =
857$, $N_2^{*} = (800 - 600)/0.7 = 286$; no equilíbrio, a espécie
1 está abaixo de seu $K$ em $\alpha N_2^{*} = 143$ e a espécie 2 em $\beta
N_1^{*} = 514$.
\end{solution}

\begin{solution}{exo:b2:speciation:6}
$\alpha K_2 = 1200 > K_1$: a espécie 1 não pode invadir a espécie 2 em
sua capacidade de suporte; $\beta K_1 = 300 < K_2$: a espécie 2 pode
invadir a espécie 1. A espécie 2 sempre vence, porque seus indivíduos
prejudicam a espécie 1 mais do que os próprios indivíduos da espécie 1,
enquanto ela mesma quase não é afetada. Com $\alpha = \beta = 1.5$: $\alpha K_2 = 1200 > 1000$
e $\beta K_1 = 1500 > 800$: nenhuma pode invadir a outra, cada uma
exclui a outra quando estabelecida, e vence a que chegou primeiro.
\end{solution}

\begin{solution}{exo:b2:speciation:7}
O tetraploide $4n$ produz gametas $2n$; cruzado com os gametas $n$ do
\omterm{def:b2:meiosis-heredity:meiosis}{diploide}, ele dá triploides $3n$, cuja \omterm{def:b2:meiosis-heredity:meiosis}{meiose} não consegue parear três
lotes e produz gametas aneuploides, na maioria inviáveis. O
tetraploide só pode se reproduzir consigo mesmo (ou por
autofecundação): isolado de uma só vez. Em geral ele falha porque está
sozinho --- seus únicos parceiros potenciais são \omterm{def:b2:meiosis-heredity:meiosis}{diploides}, e todo
cruzamento desse tipo é desperdiçado --- e é diluído; ele persiste
quando se autofecunda, se propaga vegetativamente ou surge repetidas
vezes.
\end{solution}

\begin{solution}{exo:b2:speciation:8}
$0.01\times 10\times 10 = 1$ incompatibilidade; com 30 em cada uma, $0.01\times
900 = 9$. O isolamento cresce com o quadrado do tempo de divergência:
duas populações que divergiram durante um tempo três vezes maior são
nove vezes mais isoladas --- a bola de neve.
\end{solution}

\begin{solution}{exo:b2:speciation:9}
$w = 2\sqrt{8/0.05} = \qty{25}{km}$; $w = 0.3\sqrt{8/0.4} =
\qty{1.3}{km}$. O reforço é mais provável na segunda: os híbridos são
fortemente inaptos, de modo que uma fêmea que recusa a outra forma
ganha muito, e a zona é estreita o bastante para que as duas formas se
encontrem com frequência suficiente para que essa recusa seja
selecionada.
\end{solution}

\begin{solution}{exo:b2:speciation:10}
$\alpha K_2 = 450 < 500$, de modo que cada uma invade a outra: elas
coexistem, com $N^{*} = (500 - 450)/(1 - 0.81) = 263$ cada --- por pouco, perto do
limite da exclusão, e uma pequena mudança de $K$ desfaria o equilíbrio. Com
$\alpha = \beta = 0.4$: $N^{*} = (500 - 200)/(1 - 0.16) = 357$ cada.
O deslocamento girou cada isóclina para longe da outra ($K/\alpha$
se deslocou para fora), de modo que o cruzamento se afastou dos eixos e
cada espécie fica mais perto de seu próprio $K$: uma coexistência mais
segura, com mais indivíduos de cada uma.
\end{solution}

\begin{solution}{exo:b2:speciation:11}
$R = h^{2}S = 0.7\times(-1) = \qty{-0.7}{mm}$: como observado. Na
seca de 1977, sem nenhum tentilhão-terrestre-grande presente, os mesmos
tentilhões médios foram selecionados para bicos \emph{maiores}, pois as
\omterm{def:b2:angiosperm-reproduction:seed}{sementes} grandes eram o alimento que restava; em 2004, as \omterm{def:b2:angiosperm-reproduction:seed}{sementes}
grandes foram tomadas pelo invasor, e os tentilhões médios que competiam
por elas perderam: o sentido da seleção foi fixado pelo competidor.
\end{solution}

\begin{solution}{exo:b2:speciation:12}
Feitas por acidente: o modelo de dois lócus mostra o isolamento surgindo
como subproduto de substituições fixadas por outras razões, ou por
deriva, e a geografia --- uma barreira, uma fundação --- fornece a
separação que deixa o acidente se acumular. Mantidas pela seleção: onde
as formas se encontram, a seleção contra os híbridos inaptos constrói
barreiras pré-zigóticas (reforço), e a competição desloca os caracteres
para que as duas possam coexistir. O slogan omite a \omterm{prop:b2:speciation:modes}{especiação simpátrica}
por seleção disruptiva, em que a seleção faz a espécie além de mantê-la,
e a \omterm{def:b2:genome-diversification:polyploidy}{poliploidia}, em que o acidente é a história inteira.
\end{solution}

\begin{solution}{pb:b2:speciation:1}
\textbf{1.} Espécie 1: $N_1 + 0.9N_2 = 1200$, interceptos $N_1 = 1200$
e $N_2 = 1333$. Espécie 2: $N_2 + 0.3N_1 = 400$, interceptos $N_2 =
400$ e $N_1 = 1333$.
\textbf{2.} $\beta K_1 = 360 < K_2 = 400$: a espécie 2 invade;
$\alpha K_2 = 360 < K_1 = 1200$: a espécie 1 invade. Coexistência
estável.
\textbf{3.} $N_1^{*} = (1200 - 360)/(1 - 0.27) = 1151$; $N_2^{*} =
(400 - 360)/0.73 = 55$.
\textbf{4.} $\mathrm{d}N_1/\mathrm{d}t = 0.5\times 1200\,(1 - 1218/1200)
= -9$ por ano: o residente, em sua capacidade de suporte, é empurrado
para o declínio pelos recém-chegados.
\textbf{5.} $K_1 = 600$, $K_2 = 200$: $N_1^{*} = (600 - 180)/0.73 =
575$, $N_2^{*} = (200 - 180)/0.73 = 27$: ambas caem à metade, e o
invasor, com 27 aves, fica perto da extinção --- é na seca que a
competição aperta.
\textbf{6.} O invasor é a ave maior, de bico maior: ele pega as
\omterm{def:b2:angiosperm-reproduction:seed}{sementes} grandes que o residente também usa, e as pega mais depressa,
ao passo que as \omterm{def:b2:angiosperm-reproduction:seed}{sementes} pequenas do residente pouco lhe servem.
\textbf{7.} $S = 9.0 - 9.5 = \qty{-0.5}{mm}$; $R = 0.7\times(-0.5) =
\qty{-0.35}{mm}$.
\textbf{8.} $9.5 - 0.35 = \qty{9.15}{mm}$.
\textbf{9.} $N_1^{*} = (1200 - 200)/(1 - 0.15) = 1176$; $N_2^{*} =
(400 - 360)/0.85 = 47$.
\textbf{10.} O intercepto em $N_2$ da isóclina da espécie 1 subiu de
$K_1/0.9 = 1333$ para $K_1/0.5 = 2400$: a reta girou para longe da do
invasor, o residente agora é menos afetado por cada invasor, e o
equilíbrio se deslocou em direção ao $K$ do residente.
\textbf{11.} $2/0.35 \approx 6$ anos. Isso não acontece porque uma
seleção desse tipo só ocorre em anos de seca, se inverte nos anos
chuvosos, quando as \omterm{def:b2:angiosperm-reproduction:seed}{sementes} pequenas abundam (ou quando o invasor é
escasso), e porque a variância aditiva se esgotaria.
\textbf{12.} Dois milhões de anos são cerca de $10^{6}$ gerações; um
deslocamento de \qty{0.35}{mm} por geração, mantido durante apenas um
milésimo delas, mudaria o bico em centenas de milímetros. A seleção é
forte, mas inconstante; a mudança direcional sustentada é a rara
exceção, e a mudança líquida ao longo do tempo geológico é um pequeno
resíduo de grandes episódios opostos.
\textbf{13.} O canto, aprendido com o pai e usado pelas fêmeas na
escolha do parceiro, com a própria forma do bico como sinal. Os híbridos
de bico intermediário não lidam bem nem com as \omterm{def:b2:angiosperm-reproduction:seed}{sementes} grandes nem com
as pequenas e são os primeiros a morrer de fome nas secas.
\textbf{14.} $w = 0.5\sqrt{8/0.2} = \qty{3.2}{km}$.
\textbf{15.} A ilha é mais estreita do que seria a zona: nenhuma clina
espacial pode se formar; a hibridização, onde ocorre, afeta a população
inteira, e só o cruzamento preferencial mantém as duas
separadas.
\textbf{16.} As fêmeas que só se acasalam com machos de seu próprio
canto e bico deixam mais netos que as que hibridizam; a preferência e o
sinal (canto, bico) se tornam mais distintos onde as duas espécies vivem
juntas do que onde vivem sozinhas.
\textbf{17.} $0.005\times 15\times 15 = 1.1$ incompatibilidades: em
média uma, com híbridos provavelmente de \omterm{def:b2:population-genetics:fitness}{aptidão} reduzida, mas não
estéreis. Ainda não são espécies pelo critério biológico; estão a
caminho.
\textbf{18.} Biológico: ainda não (híbridos quase aptos).
Morfológico: talvez, se os bicos ou a plumagem divergiram
visivelmente. Filogenético: sim, se cada população tiver uma diferença
diagnóstica fixada. Os conceitos discordam justamente durante o
processo de especiação, o que os torna úteis.
\textbf{19.} $1/(5\times 10^{6}) = 2\times 10^{-7}$ por espécie por
ano; para $10^{7}$ espécies, duas extinções por ano.
\textbf{20.} De 200 a 2000 espécies por ano.
\textbf{21.} $0.75\times 10^{7}/10^{5} = 75$ espécies por ano --- menos
que a taxa atual, mesmo em sua estimativa mais baixa. O episódio atual,
se mantido, é uma \omterm{prop:b2:speciation:tempo}{extinção em massa} que avança mais depressa que as cinco
grandes.
\textbf{22.} Os nichos deixados vagos pelos perdedores estão vazios de
competidores, de modo que qualquer variante sobrevivente capaz de usar
um deles é favorecida, e os sobreviventes se diversificam para ocupá-los.
\textbf{23.} \omterm{prop:b2:speciation:tempo}{Equilíbrio pontuado}; \omterm{prop:b2:speciation:modes}{especiação alopátrica} (peripátrica)
numa pequena população isolada cuja descendente depois se expande pela
área de distribuição da população principal.
\textbf{24.} $1.04^{n} = 2$: $n = \ln 2/\ln 1.04 = 18$ gerações.
O registro mostra milhões de anos porque a seleção se inverte, a \omterm{prop:b2:population-genetics:kinds}{seleção
estabilizadora} domina a maior parte do tempo, e a mudança líquida é a
pequena diferença entre grandes episódios opostos.
\textbf{25.} Antes do deslocamento: coexistência estável com $N_1^{*} =
1151$, $N_2^{*} = 55$; depois: 1176 e 47. Deslocamento do bico de
\qty{-0.35}{mm} em uma geração. Largura da \omterm{prop:b2:speciation:reinforcement}{zona híbrida} de \qty{3.2}{km}.
\end{solution}
